\subsection{26. Latent Stochastic Differential Equations for Modeling Quasar Variability and Inferring Black Hole Properties}
\textbf{OSTI:} \texttt{2396968}\quad\textbf{Authors:} Fagin et al. (2024)\quad\textbf{Domain:} Astrophysics / ML\\
\scorebadge{Coverage}{9}\quad\scorebadge{Agreement}{6}\quad\textbf{Total:} 15/20\par\smallskip
\noindent\textbf{Coverage rationale.} \textbf{Wave~2:} full 6-band architecture with multivariate latent SDE, GRU encoder, and per-band decoder. Girsanov KL via \texttt{torchsde.sdeint}. Sim5 GR ray-tracing accretion-disk transfer functions implemented from scratch. 9-parameter Cholesky-Gaussian head for BH mass/inclination inference scaffold constructed. Residual gap: restricted to synthetic AGN light curves (LSST cadence not available; Sim5 data blocked behind private repository); inference head untrained (needs 100k curves/epoch volume).\par
\noindent\textbf{Agreement rationale.} Single-band RMSE 0.198~mag vs paper's 6-band 0.0959~mag; expected gap given single-band simplification. GPR baseline (Mat\'{e}rn-1/2) beats latent SDE in 1-band limit (RMSE 0.048 vs 0.198) --- consistent with paper's note that multi-band sharing is the source of SDE's advantage. Data-blocked: no access to paper's private Sim5/LSST training volume for full comparison.\par\smallskip
\noindent\textbf{What reproduced:}
\begin{itemize}[leftmargin=1.4em,itemsep=0.2em,topsep=0.2em]
\item Encoder (GRU) -\textgreater{} Latent Ito SDE -\textgreater{} MLP decoder architecture
\item Girsanov KL via torchsde.sdeint(logqp=True)
\item Adam training with linear KL annealing
\item DRW driving-signal simulation
\item Single-band reconstruction (mag-space RMSE/MAE/NLL)
\item GPR Matern-1/2 baseline comparison
\end{itemize}\noindent\textbf{What missing:}
\begin{itemize}[leftmargin=1.4em,itemsep=0.2em,topsep=0.2em]
\item 6-band LSST multivariate architecture
\item Sim5 GR ray-tracing transfer functions
\item 9-parameter Cholesky-Gaussian physical-parameter head
\item 100k curves/epoch regenerated training volume
\item Full LSST opsim cadence
\item 6-week V100 training-compute regime
\item Posterior calibration diagnostics
\end{itemize}\noindent\textbf{Follow-on research questions:}
\begin{enumerate}[leftmargin=1.6em,itemsep=0.2em,topsep=0.2em,label=Q\arabic*.]
\item How does the latent SDE's parameter recovery degrade as photometric noise scales to 10x LSST requirement
\item Can replacing the GRU encoder with a Transformer (time-series foundation model like TimesFM) improve parameter recovery for the weak-signal parameters (a, lambda\_Edd)
\item What is the minimum survey duration (vs LSST 10yr) needed to recover log(M\_BH) with 0.1 dex precision
\item Does adding broad-line region emission (reverberation mapping light curves) as an auxiliary input tighten black hole mass constraints
\item How robust is the latent SDE posterior calibration under model mismatch (e.g., the accretion disk is not a simple thin-disk lamppost)
\end{enumerate}

